\documentclass[12pt,a4paper]{article}

\usepackage[spanish,es-nodecimaldot]{babel}
\usepackage[utf8]{inputenc}
\usepackage[T1]{fontenc}
\usepackage{lmodern}
\usepackage{geometry}
\usepackage{microtype}
\usepackage{setspace}
\usepackage{parskip}
\usepackage{enumitem}
\usepackage{xcolor}
\usepackage{listings}
\usepackage{hyperref}
\usepackage{amsmath}
\usepackage{amssymb}
\usepackage{array}
\usepackage{booktabs}
\usepackage{longtable}

\geometry{margin=2.5cm}
\onehalfspacing
\setlist[itemize]{leftmargin=*,itemsep=0.2em}
\setlist[enumerate]{leftmargin=*,itemsep=0.2em}

\hypersetup{
  colorlinks=true,
  linkcolor=black,
  urlcolor=blue,
  pdftitle={Resolucion del Practico 9 (Trabajo Practico Final) - Testing de Software 2024},
  pdfauthor={}
}

\lstdefinestyle{codestyle}{
  basicstyle=\ttfamily\small,
  frame=single,
  breaklines=true,
  showstringspaces=false,
  keywordstyle=\color{blue!70!black}\bfseries,
  commentstyle=\color{gray!70!black}\itshape,
  stringstyle=\color{orange!70!black}
}

\title{\textbf{Trabajo Practico Final --- Practico 9}\\Testing de Software 2024}
\author{}
\date{}

\begin{document}

\maketitle
\tableofcontents
\newpage

\section*{Introduccion}
\addcontentsline{toc}{section}{Introduccion}

Este documento contiene la resolucion del \textit{Trabajo Practico Final} (Practico 9) de
Testing de Software. El practico integra los conceptos vistos a lo largo del cuatrimestre
sobre una implementacion (parcial) del juego \emph{Threes} (\url{http://threesjs.com}).

\textbf{Clases provistas.}

\begin{itemize}
  \item \texttt{ThreesBoard.java}: modela el tablero de $4 \times 4$ y determina que
  movimientos estan habilitados en un estado dado.
  \item \texttt{ThreesTile.java}: modela una pieza individual del tablero. Una pieza es
  ``libre'' (valor $0$) o tiene un valor valido: $1$, $2$ o $3 \cdot 2^i$ con $i \in \mathbb{N}$
  (es decir, $3, 6, 12, 24, 48, \dots$).
  \item \texttt{ThreesController.java}: realiza los movimientos del juego
  (\texttt{moveUp}, \texttt{moveDown}, \texttt{moveLeft}, \texttt{moveRight}), gestiona la
  carga de la siguiente baldosa en el borde opuesto al movimiento y guarda
  \texttt{nextTileValue}.
  \item \texttt{RandomGenerator.java} / \texttt{RandomValueGenerator.java}: interfaz y
  implementacion concreta para la fuente de aleatoriedad. Se aprovecha en los tests para
  inyectar una secuencia determinista (\texttt{SequenceRandomGenerator}).
\end{itemize}

\textbf{Reglas del juego (resumen).}

\begin{itemize}
  \item Tablero $4 \times 4$. Cada juego inicia con $9$ piezas aleatorias con valores
  $\in \{1,2,3\}$.
  \item Cada movimiento desplaza \emph{todas} las piezas en la misma direccion a lo sumo un
  lugar.
  \item \textbf{Combinaciones:} $1+2 \to 3$ (y $2+1 \to 3$); $x + x \to 2x$ para $x \ge 3$.
  Las combinaciones $3+6$, $6+12$, etc.\ no son validas.
  \item Por fila/columna se realiza \emph{a lo sumo una} combinacion.
  \item Despues de un movimiento, una nueva baldosa (valor $1$, $2$ o $3$) entra en alguna
  posicion del borde opuesto a las filas/columnas modificadas.
  \item El juego termina cuando no hay movimiento posible.
  \item \textbf{Invariante:} cada pieza con valor $n$ distinto de $1$ y $2$ cumple
  $\exists\, i \in \mathbb{N} : n = 3 \cdot 2^i$.
\end{itemize}

\textbf{Estructura del documento.}

\begin{itemize}
  \item \textbf{Ejercicio 1.a}: completar la implementacion, reportar y corregir defectos.
  \item \textbf{Ejercicio 1.b}: suite de tests para \texttt{ThreesBoard} usando
  \emph{Caracterizacion del Dominio de Inputs}, con sus requisitos y la suite que los cubre.
  \item \textbf{Ejercicio 2}: suite de tests de modulo/integracion para \texttt{moveUp},
  \texttt{moveDown}, \texttt{moveLeft} y \texttt{moveRight} con alta cobertura de predicados
  medida con JaCoCo.
\end{itemize}

\textbf{Ejecucion de la suite y de JaCoCo.} Desde \texttt{tp9/assignment-9-rodeghiero-1}:

\begin{lstlisting}[style=codestyle]
JAVA_HOME=$(/usr/libexec/java_home -v 17) mvn -q test
\end{lstlisting}

\section{Ejercicio 1.a: completar la implementacion}

\subsection{Errores detectados y corregidos}

Durante la implementacion se encontraron varios defectos en el codigo provisto. Cada uno se
reporta a continuacion con su correccion.

\subsubsection*{(1) \texttt{ThreesBoard.getTile}: validacion incorrecta de la columna}

\textbf{Defecto.} La condicion original permitia el indice $\texttt{col} = \texttt{COLUMNS}$,
fuera de rango (\texttt{col <= COLUMNS}).

\textbf{Correccion.} Se ajustaron las dos comparaciones a estrictas en el extremo superior:

\begin{lstlisting}[style=codestyle,language=Java]
public ThreesTile getTile(int row, int col) {
    if (0 <= row && row < ROWS && 0 <= col && col < COLUMNS) {
        return elements[row][col];
    } else {
        throw new IllegalArgumentException();
    }
}
\end{lstlisting}

\subsubsection*{(2) \texttt{ThreesBoard(int totalTiles)}: sin validacion de rango}

\textbf{Defecto.} El constructor no validaba \texttt{totalTiles}. Para valores $> 16$ el
\texttt{while} interno entraba en bucle infinito (no quedaban celdas libres).

\textbf{Correccion.}

\begin{lstlisting}[style=codestyle,language=Java]
public ThreesBoard(int totalTiles) {
    if (totalTiles < 0 || totalTiles > ROWS * COLUMNS) {
        throw new IllegalArgumentException("totalTiles must be in [0, 16]");
    }
    // ... resto de la inicializacion ...
}
\end{lstlisting}

\subsubsection*{(3) \texttt{ThreesBoard}: valor aleatorio inicial incorrecto}

\textbf{Defecto.} La inicializacion aleatoria seteaba valores entre $0$ y $2$, incluyendo el
$0$ (celda libre), cuando deberia setear $1$, $2$ o $3$.

\textbf{Correccion.} Se cambia por \texttt{generator.getRandom(3) + 1} (rango $[1,3]$).

\subsubsection*{(4) \texttt{ThreesBoard}: metodos faltantes}

Estaban sin implementar:

\begin{itemize}
  \item \texttt{isFinished()}
  \item \texttt{canMoveLeft()}, \texttt{canMoveRight()}, \texttt{canMoveUp()},
  \texttt{canMoveDown()}
  \item \texttt{canTilesCombine(ThreesTile, ThreesTile)}
  \item \texttt{computeScore()}
\end{itemize}

\textbf{Decision de diseno.} \texttt{isFinished()} es simplemente el \texttt{!canMoveLeft \&\&
!canMoveRight \&\& !canMoveUp \&\& !canMoveDown}. Cada uno de los \texttt{canMove*} recorre
el tablero en su direccion y devuelve \texttt{true} si encuentra al menos una celda no libre
que pueda avanzar (porque la vecina en la direccion del movimiento esta libre o es
combinable). Se factorizo la logica de combinaciones en \texttt{canTilesCombine}.

\begin{lstlisting}[style=codestyle,language=Java]
public boolean canMoveLeft() {
    for (int i = 0; i < ROWS; i++) {
        for (int j = 1; j < COLUMNS; j++) {
            if (!elements[i][j].isFree()) {
                ThreesTile current = elements[i][j];
                ThreesTile left    = elements[i][j - 1];
                if (left.isFree() || canTilesCombine(left, current)) {
                    return true;
                }
            }
        }
    }
    return false;
}

public boolean canTilesCombine(ThreesTile t1, ThreesTile t2) {
    if (t1 == null || t2 == null) {
        throw new IllegalArgumentException("Tiles cannot be null");
    }
    if (t1.isFree() || t2.isFree()) return false;
    int v1 = t1.getValue(), v2 = t2.getValue();
    return (v1 == 1 && v2 == 2)
        || (v1 == 2 && v2 == 1)
        || (v1 >= 3 && v1 == v2);
}
\end{lstlisting}

\texttt{canMoveRight}, \texttt{canMoveUp} y \texttt{canMoveDown} son simetricos.

\textbf{\texttt{computeScore()}.}

\begin{lstlisting}[style=codestyle,language=Java]
public int computeScore() {
    int score = 0;
    for (int i = 0; i < ROWS; i++) {
        for (int j = 0; j < COLUMNS; j++) {
            int v = elements[i][j].getValue();
            if (v == 3)      score += 3;
            else if (v > 3)  score += (v - 3) * 3;
        }
    }
    return score;
}
\end{lstlisting}

\subsubsection*{(5) \texttt{ThreesTile.isValidValue}: validacion correcta del invariante}

\textbf{Implementacion.} Un valor $v$ es valido sii: $v \in \{0, 1, 2\}$ o $v = 3 \cdot 2^i$
para algun $i \ge 0$. La verificacion se reduce a chequear que $v/3$ sea potencia de dos:

\begin{lstlisting}[style=codestyle,language=Java]
public boolean isValidValue(int v) {
    if (v == 0 || v == 1 || v == 2) return true;
    if (v < 3 || v % 3 != 0)        return false;
    int quotient = v / 3;
    return isPowerOfTwo(quotient);
}

private boolean isPowerOfTwo(int n) {
    return n > 0 && (n & (n - 1)) == 0;
}
\end{lstlisting}

El truco \texttt{n \& (n-1) == 0} es la forma clasica de chequear potencia de dos en una
operacion.

\subsubsection*{(6) \texttt{ThreesTile.setRandomValue}: distribucion incorrecta}

\textbf{Defecto.} Originalmente \texttt{value = randomGenerator.nextInt() \% 3} podia
producir valores fuera de $\{1, 2, 3\}$ (incluido el $0$ o negativos para enteros aleatorios
arbitrarios).

\textbf{Correccion.} \texttt{value = randomGenerator.nextInt(3) + 1}.

\subsubsection*{(7) \texttt{ThreesController}: \texttt{moveUp}, \texttt{moveLeft},
\texttt{moveRight} faltantes y \texttt{moveDown} con defecto}

\textbf{Defecto.} \texttt{moveUp}, \texttt{moveLeft} y \texttt{moveRight} no estaban
implementados. \texttt{moveDown}, segun la version inicial, podia duplicar valores al
desplazar (no respetaba la regla de \emph{una sola} combinacion por fila/columna).

\textbf{Diseno de la correccion.} Se factorizo el algoritmo de movimiento en helpers
privados, todos documentados con JavaDoc:

\begin{itemize}
  \item \texttt{readRow}/\texttt{readColumn}: leen una linea en el orden del movimiento.
  \item \texttt{collapseAndMergeLine}: colapsa los huecos y aplica una sola combinacion
  por par adyacente.
  \item \texttt{sameLine}: compara dos lineas para saber si hubo cambios.
  \item \texttt{writeRow}/\texttt{writeColumn}: escriben de vuelta en coordenadas del
  tablero.
  \item \texttt{loadNextTileOnColumns}/\texttt{loadNextTileOnRows}: cargan la nueva baldosa
  en el borde opuesto a las filas/columnas modificadas.
\end{itemize}

\begin{lstlisting}[style=codestyle,language=Java]
private int[] collapseAndMergeLine(int[] line) {
    int[] compact = new int[line.length];
    int n = 0;
    for (int v : line) if (v != 0) compact[n++] = v;

    int[] merged = new int[line.length];
    int m = 0, i = 0;
    while (i < n) {
        if (i + 1 < n && canValuesCombine(compact[i], compact[i + 1])) {
            merged[m++] = compact[i] + compact[i + 1];
            i += 2;
        } else {
            merged[m++] = compact[i];
            i++;
        }
    }
    return merged;
}
\end{lstlisting}

\textbf{Por que asi.} Recorrer la linea ya colapsada de izquierda a derecha y avanzar de
\emph{dos en dos} cuando hay combinacion garantiza la regla ``una sola combinacion por
fila/columna'': una vez combinados dos elementos, el indice salta al siguiente par.

\subsubsection*{(8) \texttt{pom.xml}: incompatibilidad de JaCoCo con JDK reciente}

\textbf{Defecto.} \texttt{jacoco-maven-plugin 0.8.2} fallaba al instrumentar bytecode con
JDK 17 (clases con versiones superiores no soportadas).

\textbf{Correccion.} Se actualizo el plugin a \texttt{0.8.12}.

\subsection{Convenciones de codigo aplicadas}

\begin{itemize}
  \item Cada metodo agregado lleva \texttt{JavaDoc} explicito (parametros, retornos,
  excepciones).
  \item Las validaciones se hacen al inicio del metodo (\emph{fail-fast}) y lanzan
  \texttt{IllegalArgumentException} con mensaje cuando corresponde.
  \item Las constantes del tablero estan en \texttt{ThreesBoard.ROWS} y
  \texttt{ThreesBoard.COLUMNS} (no \emph{magic numbers}).
  \item La logica de aleatoriedad se concentra en \texttt{RandomGenerator}, lo que permite
  inyectar generadores deterministas en tests.
\end{itemize}

\section{Ejercicio 1.b: suite de tests para \texttt{ThreesBoard}}

\subsection{Modelo elegido}

Se utilizo \textbf{Caracterizacion del Dominio de Inputs} para modelar el espacio de pruebas
de \texttt{ThreesBoard}, ignorando \texttt{toString()} y \texttt{getRandom()} como permite el
enunciado.

\subsection{Caracteristicas y particiones}

\begin{center}
\begin{tabular}{lll}
\toprule
ID & Caracteristica & Particiones (bloques) \\
\midrule
C1 & coordenadas (\texttt{setTile}/\texttt{getTile}) & validas $[0..3]$; invalidas ($< 0$ o $\ge 4$) \\
C2 & cantidad inicial (\texttt{ThreesBoard(totalTiles)}) & validas $[0..16]$; invalidas ($< 0$ o $> 16$) \\
C3 & ocupacion (\texttt{numberOfSetTiles}) & vacio; parcial; completo \\
C4 & tipo de pareja (\texttt{canTilesCombine}) & $1+2$; $2+1$; $x+x$ con $x \ge 3$; no combinables; \texttt{null} \\
C5 & movilidad (\texttt{canMove*}, \texttt{isFinished}) & sin movs; por huecos; por combinacion; lleno bloqueado \\
C6 & categoria de valor (\texttt{computeScore}) & $\{0,1,2\}$; $3$; $> 3$ \\
\bottomrule
\end{tabular}
\end{center}

\subsection{Criterio de cobertura}

Se utiliza \textbf{Each Choice Coverage} extendido a representantes razonables de cada bloque
en cada caracteristica. En las caracteristicas con interaccion (e.g.\ canMove en cada
direccion) se ejercita cada particion en cada direccion.

\subsection{Requisitos de test}

\begin{itemize}
  \item \textbf{R1.} \texttt{setTile/getTile} dentro del rango actualizan y recuperan el
  valor de la celda. \emph{(C1 bloque valido.)}
  \item \textbf{R2.} \texttt{setTile} rechaza coordenadas fuera de rango. \emph{(C1 bloque
  invalido.)}
  \item \textbf{R3.} \texttt{getTile} rechaza coordenadas fuera de rango. \emph{(C1 bloque
  invalido.)}
  \item \textbf{R4.} \texttt{ThreesBoard(totalTiles)} con valor valido setea exactamente esa
  cantidad de celdas con valores permitidos. \emph{(C2 valido + C3 parcial.)}
  \item \textbf{R5.} \texttt{ThreesBoard(totalTiles)} rechaza \texttt{totalTiles} invalido.
  \emph{(C2 invalido.)}
  \item \textbf{R6.} \texttt{numberOfSetTiles} cuenta exactamente las celdas no vacias.
  \emph{(C3.)}
  \item \textbf{R7.} \texttt{canTilesCombine} respeta las reglas para los casos validos e
  invalidos, y rechaza \texttt{null}. \emph{(C4.)}
  \item \textbf{R8.} Tablero vacio: ningun movimiento es posible y esta finalizado.
  \emph{(C5 ``sin movs''.)}
  \item \textbf{R9.} Tablero con huecos: el movimiento por desplazamiento esta detectado en
  cada direccion. \emph{(C5 ``por huecos''.)}
  \item \textbf{R10.} Tablero con piezas combinables: cada \texttt{canMove*} detecta la
  combinacion. \emph{(C5 ``por combinacion''.)}
  \item \textbf{R11.} Tablero lleno sin combinaciones: queda finalizado. \emph{(C5 ``lleno
  bloqueado'' + C3 ``completo''.)}
  \item \textbf{R12.} \texttt{computeScore} suma $3$ para valor $3$ y $(v-3)\cdot 3$ para
  $v > 3$. \emph{(C6.)}
\end{itemize}

\subsection{Suite construida}

\textbf{Nombre:} \texttt{ThreesBoardInputDomainCharacterizationTest}.

Convenciones:

\begin{itemize}
  \item \texttt{@BeforeEach} crea un tablero limpio para garantizar aislamiento.
  \item Los rechazos parametrizados (\texttt{R2}, \texttt{R3}, \texttt{R5}) usan
  \texttt{@ParameterizedTest} con \texttt{@CsvSource}.
  \item Los casos de combinacion (\texttt{R7}) tambien son parametrizados.
  \item Los nombres de los tests son auto-explicativos
  (\texttt{setTileRejectsOutOfRangeCoordinates}, \texttt{computeScoreFollowsImplementedRules}\ldots).
\end{itemize}

\textbf{Mapeo requisito $\to$ test.}

\begin{center}
\begin{tabular}{ll}
\toprule
Requisito & Test que lo cubre \\
\midrule
R1  & \texttt{setAndGetTileWithinBounds} \\
R2  & \texttt{setTileRejectsOutOfRangeCoordinates} (parametrizado) \\
R3  & \texttt{getTileRejectsOutOfRangeCoordinates} (parametrizado) \\
R4  & \texttt{constructorWithTotalTilesSetsExactAmount} \\
R5  & \texttt{constructorRejectsInvalidTotalTiles} (parametrizado) \\
R6  & \texttt{numberOfSetTilesCountsNonFreeTiles} \\
R7  & \texttt{canTilesCombineMatchesRules} + \texttt{canTilesCombineRejectsNullTiles} \\
R8  & \texttt{emptyBoardHasNoMovesAndIsFinished} \\
R9  & \texttt{movementByGapsIsDetectedPerDirection} \\
R10 & \texttt{movementByCombinationIsDetectedPerDirection} \\
R11 & \texttt{fullBoardWithoutCombinationsIsFinished} \\
R12 & \texttt{computeScoreFollowsImplementedRules} \\
\bottomrule
\end{tabular}
\end{center}

\subsection{Requisitos inviables}

\begin{itemize}
  \item \textbf{RI1.} Verificar las \emph{posiciones exactas} de las fichas seteadas por
  \texttt{ThreesBoard(totalTiles)} sin inyectar/mokear \texttt{RandomGenerator}.
  \textbf{Motivo:} la posicion depende del generador aleatorio; sin un doble determinista
  el test seria flaky.
  \item \textbf{RI2.} Cobertura exhaustiva de todos los estados del tablero.
  \textbf{Motivo:} explosion combinatoria; con 16 celdas y multiples valores posibles
  por celda, el espacio de estados es inabordable. Por eso se aplica particionado por
  dominio de inputs.
  \item \textbf{RI3.} Verificar uniformidad estadistica del generador aleatorio como tests
  unitarios.
  \textbf{Motivo:} esa propiedad es de naturaleza estadistica, no se valida con asserts
  deterministas en una corrida unitaria.
\end{itemize}

\subsection{Resultados de cobertura con JaCoCo}

Sobre \texttt{ThreesBoard} (Ejercicio 1.b):

\begin{center}
\begin{tabular}{ll}
\toprule
Metrica & Valor \\
\midrule
Lineas    & $85/97 = \textbf{87.63\%}$ \\
Branches  & $108/124 = \textbf{87.10\%}$ \\
Metodos   & $12/13 = \textbf{92.31\%}$ \\
\bottomrule
\end{tabular}
\end{center}

\textbf{Justificacion de las partes no cubiertas.} La fraccion no cubierta corresponde
principalmente a:

\begin{itemize}
  \item Ramas defensivas que ya estan ejercitadas indirectamente por la suite (e.g.\
  comparaciones que se cumplen siempre con los valores del dominio elegido).
  \item \texttt{toString()}: explicitamente excluido del modelo segun el enunciado.
  \item Algunas ramas del constructor aleatorio que dependen de la fuente de
  aleatoriedad (cf.\ RI1).
\end{itemize}

\section{Ejercicio 2: tests de modulo/integracion para los movimientos}

\subsection{Objetivo}

Se busca \textbf{alta cobertura de predicados} de las acciones \texttt{moveUp},
\texttt{moveDown}, \texttt{moveLeft} y \texttt{moveRight} de \texttt{ThreesController},
medida con JaCoCo. Un test que efectua \emph{dos movimientos diferentes en secuencia} cuenta
como test de modulo.

\subsection{Modelo de estados por accion}

Cada accion de movimiento es la composicion de cuatro decisiones internas que el modelo
captura como cinco estados representativos:

\begin{itemize}
  \item \textbf{M0}: tablero \emph{no} modificable en esa direccion (no se mueve nada;
  \texttt{loadNextTile} no se invoca).
  \item \textbf{M1}: tablero modificable \emph{solo por desplazamiento} (huecos en la
  direccion).
  \item \textbf{M2}: tablero modificable con \emph{combinacion $1+2$ o $2+1$}.
  \item \textbf{M3}: tablero modificable con \emph{combinacion $x+x$} para $x \ge 3$.
  \item \textbf{M4}: carga de \emph{nueva baldosa} en el borde opuesto con la fila/columna
  elegida (pseudo)aleatoriamente.
\end{itemize}

\subsection{Inyeccion de aleatoriedad determinista}

Como las acciones de movimiento llaman a \texttt{generator.getRandom(\ldots)} para elegir
posicion de la nueva baldosa y para regenerar \texttt{nextTileValue}, se inyecta una
\texttt{SequenceRandomGenerator} (clase privada del test) que devuelve una secuencia
preconfigurada de enteros. Eso permite asserts exactos sobre la posicion final y el siguiente
valor cargado.

\begin{lstlisting}[style=codestyle,language=Java]
private static class SequenceRandomGenerator implements RandomGenerator {
    private final int[] sequence;
    private int index;
    SequenceRandomGenerator(int... sequence) { this.sequence = sequence; this.index = 0; }
    @Override public int getRandom(int bound) {
        if (bound <= 0) throw new IllegalArgumentException("bound must be positive");
        if (index >= sequence.length) return 0;
        int value = sequence[index++];
        return Math.floorMod(value, bound);
    }
}
\end{lstlisting}

\subsection{Requisitos de test}

\begin{itemize}
  \item \textbf{R2.1.} Constructor por defecto: tablero inicial con $9$ baldosas y
  \texttt{nextTileValue} valido en $[1, 3]$.
  \item \textbf{R2.2.} Si el tablero no admite cambios en una direccion, el movimiento
  retorna \texttt{false} y \emph{no} cambia ni el tablero ni \texttt{nextTileValue}. Se
  parametriza en las cuatro direcciones (M0).
  \item \textbf{R2.3.} \texttt{moveLeft} combina, desplaza y carga la nueva baldosa en el
  borde derecho (M1 + M3).
  \item \textbf{R2.4.} \texttt{moveRight} con combinacion $2+1$ y carga en el borde
  izquierdo (M2).
  \item \textbf{R2.5.} \texttt{moveUp} combina y carga en la fila inferior (M3).
  \item \textbf{R2.6.} \texttt{moveDown} combina y carga en la fila superior (M3 + M4).
  \item \textbf{R2.7.} Test de modulo: dos movimientos diferentes en secuencia se comportan
  consistentemente (M2/M3 + verificacion de estado global).
\end{itemize}

\subsection{Suite construida}

\textbf{Nombre:} \texttt{ThreesControllerModuleIntegrationTest}.

\textbf{Mapeo requisito $\to$ test.}

\begin{center}
\begin{tabular}{ll}
\toprule
Requisito & Test que lo cubre \\
\midrule
R2.1 & \texttt{defaultConstructorCreatesExpectedInitialState} \\
R2.2 & \texttt{noChangeMoveKeepsBoardAndNextTile} (parametrizado, 4 direcciones) \\
R2.3 & \texttt{moveLeftCombinesShiftsAndLoadsNextTileOnRightEdge} \\
R2.4 & \texttt{moveRightCombinesTwoPlusOneAndLoadsNextTileOnLeftEdge} \\
R2.5 & \texttt{moveUpCombinesAndLoadsNextTileOnBottomRow} \\
R2.6 & \texttt{moveDownCombinesAndLoadsNextTileOnTopRow} \\
R2.7 & \texttt{sequentialDifferentMovesBehaveConsistently} \\
\bottomrule
\end{tabular}
\end{center}

\textbf{Ejemplo de test concreto.} Para \texttt{moveLeft}:

\begin{lstlisting}[style=codestyle,language=Java]
@Test
@DisplayName("moveLeft: combina, desplaza y carga nueva baldosa en borde derecho")
void moveLeftCombinesShiftsAndLoadsNextTileOnRightEdge() {
    load(board, new int[][] {
        {0, 1, 2, 0},
        {3, 3, 0, 0},
        {6, 12, 0, 0},
        {0, 0, 0, 0}
    });
    board.generator = new SequenceRandomGenerator(
        1, // constructor: nextTileValue = 2
        1, // moveLeft: elige movedRows[1] => fila 1
        0  // nuevo nextTileValue = 1
    );
    ThreesController controller = new ThreesController(board);

    boolean moved = controller.moveLeft();

    assertTrue(moved);
    assertBoardEquals(new int[][] {
        {3, 0, 0, 0},   // 1+2 = 3, queda en columna 0
        {6, 0, 0, 2},   // 3+3 = 6; baldosa nueva (valor 2) en borde derecho
        {6, 12, 0, 0},  // sin combinacion: 6 y 12 no se combinan
        {0, 0, 0, 0}
    }, board);
    assertEquals(1, controller.getNextTileValue());
}
\end{lstlisting}

\textbf{Comentarios.}

\begin{itemize}
  \item La secuencia inyectada permite asserts \emph{exactos} sobre la posicion de la
  baldosa nueva.
  \item El test ejercita simultaneamente M1 (desplazamiento del $1$ al $0$), M2
  ($1 + 2$) y M3 ($3 + 3$).
  \item La fila 2 ($6, 12, 0, 0$) demuestra que las piezas con $x \ne y$ y $x \ge 3$
  \emph{no} se combinan.
\end{itemize}

\subsection{Requisitos inviables}

\begin{itemize}
  \item \textbf{RI2.1.} Comprobar la \emph{distribucion estadistica} de la baldosa nueva
  usando aleatoriedad real. \textbf{Motivo:} no determinista; en este nivel se trabaja con
  \texttt{SequenceRandomGenerator}, dejando la propiedad estadistica fuera del scope de
  integracion reproducible.
  \item \textbf{RI2.2.} Cobertura exhaustiva de todas las secuencias posibles de estados
  del juego. \textbf{Motivo:} explosion combinatoria; se aplica particionado representativo
  del dominio (M0-M4).
\end{itemize}

\subsection{Resultados de cobertura con JaCoCo}

Sobre \texttt{ThreesController} (Ejercicio 2):

\begin{center}
\begin{tabular}{ll}
\toprule
Metrica & Valor \\
\midrule
Instrucciones & $596/596 = \textbf{100.00\%}$ \\
Branches      & $71/78 = \textbf{91.03\%}$ \\
Lineas        & $124/124 = \textbf{100.00\%}$ \\
Metodos       & $19/19 = \textbf{100.00\%}$ \\
\bottomrule
\end{tabular}
\end{center}

\textbf{Justificacion de las branches no cubiertas (7 sobre 78).} Las ramas no cubiertas
corresponden a casos defensivos como:

\begin{itemize}
  \item Comprobaciones simetricas en \texttt{loadNextTileOnColumns} y
  \texttt{loadNextTileOnRows} para movimientos sin cambios (\texttt{movedColumns} o
  \texttt{movedRows} vacios) que ya quedan cubiertos por el flujo principal.
  \item Algun camino interno del helper \texttt{canValuesCombine} que es alcanzable solo con
  combinaciones que no aparecen en las configuraciones representativas elegidas (por
  ejemplo, $a == 1 \wedge b \neq 2$ con $b = 0$).
\end{itemize}

Se opto por no agregar tests sintetizados solo para forzar cobertura de esas ramas, ya que
no aportan deteccion adicional y la consigna pide \emph{justificar} en lugar de inflar
artificialmente la suite.

\subsection{Resultado global}

\begin{lstlisting}[style=codestyle]
JAVA_HOME=$(/usr/libexec/java_home -v 17) mvn -q test
\end{lstlisting}

\begin{itemize}
  \item Total de tests ejecutados: \textbf{38} (entre ambas suites).
  \item Fallos: \textbf{0}.
  \item Errores: \textbf{0}.
\end{itemize}

\section*{Conclusiones}
\addcontentsline{toc}{section}{Conclusiones}

\begin{itemize}
  \item El trabajo cierra el ciclo completo del curso: \emph{leer codigo}, \emph{completar
  la implementacion} (con deteccion y correccion de defectos), \emph{modelar} el espacio
  de pruebas con una tecnica formal (caracterizacion del dominio de inputs) y
  \emph{verificar} la suite con cobertura medida por una herramienta externa (JaCoCo).
  \item La separacion entre \texttt{ThreesBoard} (modelo + reglas) y \texttt{ThreesController}
  (acciones) hizo natural dividir el practico en dos niveles: tests \emph{unitarios} sobre el
  tablero y tests \emph{de modulo/integracion} sobre el controlador.
  \item La inyeccion de \texttt{RandomGenerator} como dependencia permitio escribir asserts
  exactos sobre acciones que internamente son aleatorias. Sin esa inyeccion, varios
  requisitos habrian quedado en la categoria de ``inviables''.
  \item Los requisitos marcados como inviables (RI1-RI3 y RI2.1-RI2.2) son intencionalmente
  honestos: explicitan los limites de un test unitario o de integracion deterministico y
  separan lo que puede testearse de lo que requeriria otro tipo de analisis (estadistico o
  combinatorio).
\end{itemize}

\end{document}
